% Section 3.5, from lectures/L03/appendix/c_exercises.md.

\section{Exercises}
\label{c3:sec:exercises}\label{c3:app:c}

\begin{aside}
\textbf{How to work these.} These exercises reinforce the concepts from
\secref{c3:sec:inheritance} and \secref{c3:sec:interfaces}. Unlike the earlier chapters, all four
sets build \emph{one} program: an interface, two drivers implementing it, and a \code{main()} that
uses both through the interface. Build it with the Makefile from \secref{c1:sec:tips}, listing the
three source files and adding \code{-Iinclude} to the compiler flags.

The solution is in \secref{sol:sec:c3}. Try each exercise yourself before looking at it.
\end{aside}

Create the following directory structure:

\begin{console}
Makefile
include/
    driver/
        serial/
            console.hpp
            interface.hpp
            stub.hpp
source/
    driver/
        serial/
            console.cpp
            stub.cpp
    main.cpp
\end{console}

\note When a class overrides a method that is marked \code{[[nodiscard]]} in the interface, repeat
the attribute on the overriding method as well. The attribute is not inherited; see
\secref{c3:sec:structure}.

% ------------------------------------------------------------------------------------------
\exerciseset{Serial Interface}
\label{c3:set:1}

\exercise{Interface for serial communication}{Code}
\label{c3:ex:1.1}
In this exercise you will design an interface \code{driver::serial::Interface} representing a
generic serial communication driver.

\task{Tasks}
In \file{driver/serial/interface.hpp}, design a class named \code{driver::serial::Interface} that
represents a generic serial driver.

All methods except the destructor shall be purely virtual (\code{= 0}).

\task{a) Destructor}
Add a destructor that is:
\begin{itemize}
  \item Virtual.
  \item Marked \code{noexcept}.
  \item Implemented using \code{= default}.
\end{itemize}

\task{b) Initialization status}
Add a pure virtual method \code{isInitialized()} that:
\begin{itemize}
  \item Indicates whether the driver has been initialized (\code{true/false}).
  \item Does not modify the object (\code{const}).
  \item Cannot throw exceptions.
  \item Generates a warning if the return value is discarded (\code{[[nodiscard]]}).
\end{itemize}

\task{c) Transmitting data}
Add a pure virtual method \code{write()} that:
\begin{itemize}
  \item Transmits one byte of data.
  \item Takes a parameter of type \code{std::uint8_t}.
  \item Does not return a value.
  \item Cannot throw exceptions.
\end{itemize}

\task{d) Receiving data}
Add a pure virtual method \code{read()} that:
\begin{itemize}
  \item Attempts to read one byte from the serial driver.
  \item Takes a reference to a variable where the byte will be stored.
  \item Returns \code{true} if a byte was received, otherwise \code{false}.
  \item Cannot throw exceptions.
  \item Lets the user discard the return value if desired (\code{[[nodiscard]]} omitted).
\end{itemize}

% ------------------------------------------------------------------------------------------
\exerciseset{Stub Implementation}
\label{c3:set:2}

\exercise{Serial stub}{Code}
\label{c3:ex:2.1}
In \file{driver/serial/stub.hpp}, implement a stub driver \code{driver::serial::Stub} for the
serial interface.

A stub driver simulates hardware behaviour and is useful for testing without real hardware.

The class shall:
\begin{itemize}
  \item Inherit from \code{driver::serial::Interface}.
  \item Be marked \code{final}.
\end{itemize}
Method definitions shall be placed in \file{driver/serial/stub.cpp}.

\task{a) Member variables}
Add three private member variables:
\begin{itemize}
  \item The first member variable shall:
  \begin{itemize}
    \item Store the most recently transmitted byte.
    \item Have the type \code{std::uint8_t}.
    \item Be named \code{myLastByte}.
  \end{itemize}
  \item The second member variable shall:
  \begin{itemize}
    \item Indicate whether the driver is initialized.
    \item Have the type \code{bool}.
    \item Be named \code{myInitialized}.
  \end{itemize}
  \item The third member variable shall:
  \begin{itemize}
    \item Indicate whether a byte is available to read.
    \item Have the type \code{bool}.
    \item Be named \code{myHasData}.
  \end{itemize}
\end{itemize}

\task{b) Constructor}
Add a default constructor that:
\begin{itemize}
  \item Sets the driver as initialized.
  \item Indicates that no data is available to read.
\end{itemize}

\task{c) Disable copy and move semantics}
Delete the following functions (in the public section of the class):
\begin{itemize}
  \item Copy constructor.
  \item Move constructor.
  \item Copy assignment operator.
  \item Move assignment operator.
\end{itemize}

\task{d) Implement interface methods}
Implement all methods required by the interface:
\begin{itemize}
  \item The method \code{isInitialized()} shall:
  \begin{itemize}
    \item Return the value stored in \code{myInitialized}.
  \end{itemize}
  \item The method \code{write()} shall:
  \begin{itemize}
    \item Do nothing if \code{myInitialized} is \code{false}.
    \item Otherwise store the transmitted byte in \code{myLastByte}, and set \code{myHasData} to
      \code{true}.
  \end{itemize}
  \item The method \code{read()} shall:
  \begin{itemize}
    \item Return \code{false} if \code{myInitialized} is \code{false} or \code{myHasData} is
      \code{false}.
    \item Otherwise copy the stored byte (\code{myLastByte}) into the output argument, set
      \code{myHasData} to \code{false}, and return \code{true}.
  \end{itemize}
\end{itemize}

\task{e) Initialization method}
Add an additional public method \code{setInitialized()} that is not part of the interface.

This method shall:
\begin{itemize}
  \item Be used to simulate the initialization state.
  \item Take the initialization state as input.
  \item Not return a value.
\end{itemize}

% ------------------------------------------------------------------------------------------
\exerciseset{Singleton Driver}
\label{c3:set:3}

\exercise{Console serial driver}{Code}
\label{c3:ex:3.1}
In \file{driver/serial/console.hpp}, create a real driver implementation
\code{driver::serial::Console} that writes transmitted bytes to the system console.

This driver shall follow the so-called singleton pattern, i.e.\ there is only one instance of the
class.

The reason is that the system only has one console output device.

This class shall:
\begin{itemize}
  \item Inherit from \code{driver::serial::Interface}.
  \item Be marked \code{final}.
\end{itemize}
Method definitions shall be placed in \file{driver/serial/console.cpp}.

\task{a) Private constructor}
Add a private default constructor that prevents users from directly creating objects of the class.

\task{b) Private destructor}
Add a private destructor, mark it \code{noexcept} and \code{default}.

\task{c) Prevent multiple instances}
Ensure that only one instance of the class can exist.

Delete the following functions (in the public section of the class):
\begin{itemize}
  \item Copy constructor.
  \item Move constructor.
  \item Copy assignment operator.
  \item Move assignment operator.
\end{itemize}

\task{d) Singleton access method}
Add a static public method \code{instance()} that returns a reference to the unique instance of the
class.

Requirements:
\begin{itemize}
  \item The method shall return the console driver as a reference to
    \code{driver::serial::Interface} (even though the instance itself is of type
    \code{driver::serial::Console}).
  \item The method cannot throw exceptions.
  \item In this method:
  \begin{itemize}
    \item Create a local static instance of the driver (\code{static driver::serial::Console}).
    \item Initialize the instance using the default constructor.
  \end{itemize}
\end{itemize}

\task{e) Implement interface methods}
Implement all methods required by the interface:
\begin{itemize}
  \item The method \code{isInitialized()} shall:
  \begin{itemize}
    \item Always return \code{true}, since the console driver is always available.
  \end{itemize}
  \item The method \code{write()} shall:
  \begin{itemize}
    \item Cast the given byte to a character and print it to the console using
      \code{std::printf()} from \code{<cstdio>}.
  \end{itemize}
  \item The method \code{read()} shall:
  \begin{itemize}
    \item Always return \code{false}, since the console does not support reading.
  \end{itemize}
\end{itemize}

% ------------------------------------------------------------------------------------------
\exerciseset{Using the Interface}
\label{c3:set:4}

\exercise{Sending a message}{Code}
\label{c3:ex:4.1}
\task{a) The \code{sendMessage()} function}
In \file{source/main.cpp}, create a function \code{sendMessage()} that sends a message using the
serial interface.

The function shall:
\begin{itemize}
  \item Take a reference to \code{driver::serial::Interface}.
  \item Transmit the string \code{"Transmitting data with a serial driver!"}.
  \item Send the message one byte at a time.
\end{itemize}

\task{b) Testing the console driver}
In function \code{main()}, test the console driver:
\begin{itemize}
  \item Obtain the singleton instance of the console driver.
  \item Pass it to the \code{sendMessage()} function.
  \item Verify that the program prints the following:
\end{itemize}

\begin{console}
Transmitting data with a serial driver!
\end{console}

\task{c) Testing the stub driver}
In function \code{main()}, test the stub driver:
\begin{itemize}
  \item Create an instance of the stub driver and pass it to the function \code{sendMessage()}.
  \item Verify that the last transmitted byte \code{!} can later be read using the \code{read()}
    method.
\end{itemize}
